\section{Evaluation methodology}

\paragraph{Frozen workloads.}
The compatibility census uses 59 workloads, each run under three execution policies, giving 177
declared cells. Workloads are frozen before measurement and split into development, validation and
held-out sets; the split is recorded per row. A tier-A workload runs unmodified; a tier-C workload
is an explicitly declared adaptation. Tier-A and tier-C cells are never summed, and no
native-policy figure is ever averaged with a controlled-contract figure.

\paragraph{Independent oracle.}
Every invocation is compared against a native execution of the same command from the same
pre-state, on standard output, standard error, exit status and a digest of the workspace tree.
Metadata is compared separately, because a content-only comparison cannot tell two fixtures
written a second apart from two genuinely identical ones. A sample counts as a hit only when the
product reports a reuse \emph{and} the cache-decision field agrees \emph{and} the observation
matches. A hit that disagrees with the oracle fails the run rather than being averaged in.

\paragraph{Paired, randomised sampling.}
Timing uses paired samples: for each repetition both arms run once from the same pre-state in
random order, and the confidence interval is computed on the per-pair ratio by percentile
bootstrap. Cold tracing needs an empty cache and therefore cannot be paired; those rows are
labelled unpaired and reported with independent intervals.

\paragraph{Session modes.}
Two modes are reported separately and never mixed. \emph{Cold} starts from a new cache, so the
first trace, the hits, the refusals and the deoptimisations are all charged to the arm --- this is
what a new user pays. \emph{Steady} populates the cache outside the timed window and reports only
the measured one. Every row states its mode.

\paragraph{Per-invocation accounting.}
The harness records each invocation's own decision, so a session's saving is the sum of decisions
that were individually verified, not an inference from the total.

\paragraph{Competitor configurations.}
\firebuild{} 0.8.2 (Ubuntu packages, user-space extraction) and 0.8.6.2 (source build, pinned
dependencies) are run with their default configuration, with a strict configuration
(\texttt{ignore-time-queries} removed), and --- for \texttt{sed}, which \firebuild{} skips by
default --- with its skip list cleared. The version and configuration are named with every number.
Statistics are collected with \texttt{-s}, which places \firebuild's own reporting inside the timed
interval; we disclose that instead of subtracting it, because removing it would remove the evidence
that the invocation reused. \texttt{ccache} 4.10.2 is measured on compilation separately.

\paragraph{Runtime versions and hosts.}
The primary timing host is one Ubuntu 24.04 machine (kernel 6.1.102, glibc 2.39-0ubuntu8.12,
ext4, 8 vCPUs). A second userspace (Ubuntu 25.10, glibc 2.42) on the same kernel was exercised for
correctness and refusal, not timing. Two GitHub Actions runners executed the project's own
correctness, smoke and environment workflows on kernels 6.1 and 7.0.0-1017-azure; the latter has
Landlock ABI~8. CI-hosted execution is a second kernel and userland, not independent third-party
reproduction, and no such reproduction exists.

\paragraph{Instrumentation.}
Timing runs the product's plain reporting mode. The JSON mode serialises the entire cached standard
output as a byte array and was measured at up to 60\,ms on a 70\,ms hit; it is used once per row,
outside the timer, to calibrate the decision channel. The plain report is written as one buffered
block after the run, and each sample's output is verified against a reference digest, so the
instrumentation cannot improve a number by skipping verification.

\section{Results}
\label{sec:results}

\begin{table}[t]
\centering
\small
\caption{Frozen compatibility census. Eligibility and reuse are counted per (tier, contract) cell and never summed across cells.}
\label{tab:compatibility}
\begin{tabular}{lrrrrrrr}
\toprule
\textbf{Tier} & \textbf{Contract} & \textbf{Declared} & \textbf{Attempted} & \textbf{Eligible} & \textbf{Reused} & \textbf{Reuse rate} & \textbf{Incorrect} \\
\midrule
A & native & 50 & 49 & 2 & 1 & 2.0 \% & 0 \\
A & allocator-entropy-v1 & 50 & 49 & 12 & 10 & 20.4 \% & 0 \\
A & compiler-runtime-v1 & 50 & 49 & 19 & 17 & 34.7 \% & 0 \\
C & native & 9 & 5 & 0 & 0 & 0.0 \% & 0 \\
C & allocator-entropy-v1 & 9 & 5 & 0 & 0 & 0.0 \% & 0 \\
C & compiler-runtime-v1 & 9 & 5 & 0 & 0 & 0.0 \% & 0 \\
**all** &  & 177 & 162 &  & 28 & 17.3 \% & 0 \\
\bottomrule
\end{tabular}
\end{table}


\paragraph{Eligibility.}
Table~\ref{tab:compatibility} is the frozen census. Native execution reaches $4.1\%$ eligibility
and $2.0\%$ reuse, because ordinary programs read entropy or the clock through libc. The audited
\texttt{allocator-entropy-v1} contract lifts eligibility to $24.5\%$ and reuse to $20.4\%$ by
supplying the allocator's entropy reproducibly for one audited glibc build. The
\texttt{compiler-runtime-v1} contract additionally supplies the clock the compiler driver reads,
reaching $38.8\%$ eligibility and $34.7\%$ reuse. Tier-C adapted workloads reach no reuse under any
policy. Across the census 28 of 162 attempted rows reused, and no row produced an incorrect
observation.

\begin{table}[t]
\centering
\small
\caption{Same-host session cost, 25 invocations, cold cache (the first trace is charged to the arm). Milliseconds; lower is better.}
\label{tab:competitive}
\begin{tabular}{lrrrrrrr}
\toprule
\textbf{Workload} & \textbf{Native} & \textbf{TraceJIT} & \textbf{Firebuild 0.8.2} & \textbf{Firebuild 0.8.6.2} & \textbf{Firebuild 0.8.6.2 tuned} & \textbf{TraceJIT steady} & \textbf{HITs (TraceJIT)} \\
\midrule
bzip2 -9 -k (1 MiB incompressible, file output) & 2846.7 & 545.3 & 2976.9 & 3119.4 & 3131.5 & 444.7 & 24 \\
gzip -9 -n -k (1 MiB, file output) & 882.7 & 425.0 & 1154.6 & 254.7 & 243.4 & 373.8 & 120 \\
gzip-mutating & 1953.9 & 1079.4 & 949.6 & 959.1 & 885.3 & 1029.2 & 16 \\
gzip -9 -n -c (1.2 MiB text, stdout) & 1856.3 & 276.9 & 299.9 & 268.0 & 260.7 & 167.9 & 120 \\
sed (200k lines) & 1824.1 & 885.8 & 2071.2 & 2113.9 & 561.4 & 606.6 & 120 \\
sha256sum (1 MiB) & 94.5 & 377.2 & 282.1 & 247.9 & 217.8 & 363.9 & 120 \\
\bottomrule
\end{tabular}
\end{table}


\paragraph{Session cost.}
Table~\ref{tab:competitive} is a same-host, cold-cache comparison at 25 invocations. The clearest
result is an applicability result, not a speed one: on \texttt{bzip2 -9 -k} both \firebuild{}
versions decline to cache at all (they refuse to cache a command that sets its output's
timestamps), so the only comparison available is against native execution, where \tracejit{}
reduces a 2846.7\,ms session to 545.3\,ms --- $5.22\times$, $80.8\%$ less cumulative time. On
\texttt{gzip} with file output, \firebuild{} 0.8.6.2 reuses and is much faster than we are. On
\texttt{gzip} to standard output the two engines are close. On \texttt{sed}, \firebuild{} is slower
by default and faster when configured to cache it. On \texttt{sha256sum} neither engine is
profitable, because a hit re-hashes exactly what the command would have hashed.

\begin{table}[t]
\centering
\small
\caption{Measured 25-invocation session economics, cold cache. The percentage is the change in cumulative execution time against native; negative values are losses.}
\label{tab:session_economics}
\begin{tabular}{lrrrrr}
\toprule
\textbf{Workload} & \textbf{Native total (ms)} & \textbf{TraceJIT total (ms)} & \textbf{Ratio} & \textbf{Cumulative change} & \textbf{HITs} \\
\midrule
bzip2 -9 -k (1 MiB incompressible, file output) & 2846.7 & 545.3 & 5.22× & +80.8 \% & 24 \\
gzip -9 -n -k (1 MiB, file output) & 882.7 & 425.0 & 2.08× & +51.9 \% & 120 \\
gzip-mutating & 1953.9 & 1079.4 & 1.81× & +44.8 \% & 16 \\
gzip -9 -n -c (1.2 MiB text, stdout) & 1856.3 & 276.9 & 6.70× & +85.1 \% & 120 \\
sed (200k lines) & 1824.1 & 885.8 & 2.06× & +51.4 \% & 120 \\
sha256sum (1 MiB) & 94.5 & 377.2 & 0.25× & -299.3 \% & 120 \\
\bottomrule
\end{tabular}
\end{table}


\paragraph{Break-even.}
Table~\ref{tab:break_even} separates four quantities that were previously conflated. The first
invocation at which cumulative cost \emph{equals} native is not the first at which it is
\emph{cheaper}: at equality the saving is zero. Model break-even, computed from one cold trace plus
steady-state hits, differs from measured break-even at sampled session lengths, and the gap is
large exactly where the hit cost is close to the native cost --- for a 1\,MiB incompressible
\texttt{gzip} the model gives three invocations while the measured session first shows a saving at
ten. We report the measured quantity in the paper and the model quantity only as a model.

\begin{table}[t]
\centering
\small
\caption{Cost across authentic artifact sizes. Cold is the first traced execution, warm is a verified HIT of the same command.}
\label{tab:artifact_scale}
\begin{tabular}{lrrrrrr}
\toprule
\textbf{Artifact} & \textbf{Workload} & \textbf{Native (ms)} & \textbf{Cold trace (ms)} & \textbf{Warm HIT (ms)} & \textbf{Native/HIT} & \textbf{Break-even n} \\
\midrule
1MiB & gzip-stdout & 38.1 & 56.3 & 3.1 & 12.1× & 2 \\
1MiB & gzip-stdout-incompressible & 27.0 & 48.2 & 4.8 & 5.6× & 2 \\
1MiB & sha256sum & 1.4 & 33.5 & 11.3 & 0.1× & never \\
1MiB & bzip2-stdout & 64.5 & 86.3 & 4.3 & 15.1× & 2 \\
1MiB & gzip-file & 42.3 & 63.3 & 8.2 & 5.1× & 2 \\
10MiB & gzip-stdout & 486.2 & 601.9 & 11.4 & 42.7× & 2 \\
10MiB & gzip-stdout-incompressible & 266.1 & 304.2 & 20.0 & 13.3× & 2 \\
10MiB & sha256sum & 8.2 & 45.4 & 11.0 & 0.7× & never \\
10MiB & bzip2-stdout & 555.8 & 733.1 & 11.9 & 46.6× & 2 \\
10MiB & gzip-file & 478.4 & 602.6 & 16.2 & 29.4× & 2 \\
\bottomrule
\end{tabular}
\end{table}


\paragraph{Artifact scale.}
Table~\ref{tab:artifact_scale} varies artifact size on a pinned release archive. Warm hits stay in
the single-digit to low-tens of milliseconds while native cost grows with the artifact, so the
ratio improves with size; the cost that grows is verified replay itself. A verified replay of a
100\,MiB artifact was measured at 79.8\,ms when the produced output is 19.9\,MB of compressed text
and 257.8\,ms when the output is 100\,MiB of incompressible data. A decomposition of the hit path
attributes roughly 39\,ms to guard validation (input hashing), 19--86\,ms to replaying the cached
stream, and up to 57\,ms to writing it. These components are measured serially, but they are not a
lower bound: a different implementation can overlap reading, hashing, copying and writing, and the
input size is not the output size. We therefore state the measured cost and the measured component
shares, and we make no claim about how fast verified replay \emph{must} be.

\begin{table}[t]
\centering
\small
\caption{Enforcement capability and the lifecycle outcome it produces. A host without Landlock ABI 3+ can never produce PROVEN; that is the documented fallback, not a defect.}
\label{tab:capability}
\begin{tabular}{lrrrrrr}
\toprule
\textbf{Environment} & \textbf{Landlock ABI} & \textbf{seccomp} & \textbf{Discovery class} & \textbf{Enforced run} & \textbf{Enforced reuse} & \textbf{Producer under --enforce} \\
\midrule
Ubuntu 24.04.5 LTS / 7.0.0-38-generic & 8 & yes & GUARDED & PROVEN & Reused & UNKNOWN \\
Ubuntu 24.04.5 / 6.1.102 & absent & yes & GUARDED & GUARDED & not demonstrated & not measured \\
ubuntu-25.10 / 6.1.102 & absent & unknown & refused & not PROVEN & not demonstrated & not measured \\
ubuntu-25.10-glibc2.42 / 6.1.102 & absent & unknown & refused & not PROVEN & not demonstrated & not measured \\
\bottomrule
\end{tabular}
\end{table}


\paragraph{Enforcement capability.}
Table~\ref{tab:capability} is the capability matrix. On a host with Landlock ABI~8 and seccomp, a
read-only command completes the full lifecycle and the reused record is \proven{}; a host with no
Landlock never labels anything \proven{} and never reuses under \texttt{--enforce}, which is the
documented fallback rather than a failure. Output producers cannot be promoted for the reason given
in Section~\ref{sec:correctness}: the policy needs the declared output to exist, and the command
recreates it with \texttt{O\_EXCL} over that existing file, producing an unmodelled failed write.

\begin{table}[t]
\centering
\small
\caption{Adversarial and negative results, including the counterexamples that are kept rather than removed.}
\label{tab:correctness}
\begin{tabular}{lrrrr}
\toprule
\textbf{Class} & \textbf{Case} & \textbf{Observed decision} & \textbf{Observation} & \textbf{Verdict} \\
\midrule
Validation-to-restore race & before\_validation & Executed & detected before exposure & expected \\
Validation-to-restore race & before\_load & Executed & detected before exposure & expected \\
Validation-to-restore race & reverted\_mutation & Executed & detected before exposure & see report \\
Validation-to-restore race & after\_validation\_before\_apply & Reused & stale bytes served & documented counterexample \\
Validation-to-restore race & during\_hashing & Executed & 8/8 concurrent-writer trials deoptimised & holds under concurrent writes \\
Cached-object corruption & one CAS object corrupted, then reuse attempted & Executed & re-executed, no hit & correct \\
Output timestamp & restored object older than the archive built from it & 2/4 projects reproduced a stale member & make treats the archive as up to date & known limitation \\
Metadata-only observation & mtime/atime/counterexample matrix & 9 cases, 0 incorrect hits & every case matches a native observation & known limitation \\
\bottomrule
\end{tabular}
\end{table}


\paragraph{Adversarial correctness.}
Table~\ref{tab:correctness} keeps the counterexamples. The post-validation mutation is the
important one: it is a real stale hit, it is reproducible, and it is excluded from the guarantee by
statement rather than by omission. The corrupted-object row is the contrast case: corrupting a
cached object costs a re-execution and produces the correct observation, because the object is
re-hashed before exposure.
